We often denote vectors by boldface letters and for a vector $\bm v = (v_1,\dots,v_d)$ we write $\bm v[i,j] := (v_i,\dots,v_j)$
for all $1 \le i \le j \le d$ and if $i=j$, we simply write $\bm v(i)$.
We also write $\bm n$ for a number $n$ to denote a vector $(n,\dots,n)$ of appropriate dimension.

An \emph{alphabet} is a finite set $\Sigma$ of \emph{symbols} (a.k.a.\ \emph{tokens}).
We write $\Sigma^*$ for the set of all \emph{words} (a.k.a.\ \emph{sequences}, \emph{strings}) of the form 
$a_1\dots a_n$, where $n \ge 0$ and $a_i \in \Sigma$ for all $i \in [1,n]$.
A \emph{language} is a subset $L \subseteq \Sigma^*$.
We assume familiarity with basic concepts in formal language theory and
complexity theory (see e.g.\ \citep{Kozen,Sipser97}). In particular, we will 
deal with finite automata. We will also use the following complexity classes:
\[
    \P \subseteq \NP \subseteq \PSPACE \subseteq \EXP \subseteq \NEXP \subseteq 
\EXPSPACE.
\]
The complexity classes $\P$ and $\NP$ correspond to problems solvable in
polynomial (resp. nondeterminsitic polynomial) time, and are well-known.
The complexity classes $\EXP$ and $\NEXP$ are similar to $\P$ and $\NP$, but
we allow the algorithm to use exponential time. The complexity classes 
$\PSPACE$ and $\EXPSPACE$ correspond to problems solvable
in polynomial (resp. exponential) space. The above inclusions are well-known
(cf. \citep{Sipser97}).


\subsection{Linear Temporal Logic}
A formula in Linear Temporal Logic (LTL) over the finite alphabet $\Sigma$ has the following syntax:
\[\varphi ::= \top \mid \bot \mid  Q_a (\text{for all } a \in \Sigma) \mid 
\varphi \wedge \varphi \mid \varphi \vee \varphi \mid \neg\varphi \mid
\varphi \since \varphi \mid \varphi \until \varphi\] 
We define satisfaction of an LTL formula $\varphi$ on a word $w = a_1\dots a_n \in \Sigma^*$ at position $i \in [1,n]$,
written $w,i \models \varphi$, inductively (omitting $\top$ (true) and $\bot$ (false)):

\begin{center}
\begin{tabular}{l l l}
$w,i \models Q_a$ & iff & $a_i = a$ \qquad (for all $a \in \Sigma$)\\
$w,i \models \varphi_1 \wedge \varphi_2$ & iff & $w,i \models \varphi_1$ and $w,i \models \varphi_2$\\
$w,i \models \varphi_1 \vee \varphi_2$ & iff & $w,i \models \varphi_1$ or $w,i \models \varphi_2$\\
$w,i \models \neg\varphi_1$ & iff & $w,i \not\models \varphi_1$\\
$w,i \models \varphi_1 \since \varphi_2$ & iff & for some $j$ with $1 \le j < i$ we have $w,j \models \varphi_2$ and \\
&& for all $k$ with $j < k < i$ we have $w,k \models \varphi_1$\\
$w,i \models \varphi_1 \until \varphi_2$ & iff & for some $j$ with $i < j \le n$ we have $w,j \models \varphi_2$ and \\
&& for all $k$ with $i < k < j$ we have $w,k \models \varphi_1$
\end{tabular}
\end{center}

Moreover, we define the shortcuts 
\[\past\varphi := \top \since \varphi \qquad \future\varphi := \top \until \varphi \qquad
\next\varphi := \bot \until \varphi \qquad \Global\varphi := \varphi \wedge \neg\future\neg\varphi.\]
An LTL formula recognizes the language $L(\varphi)$ of all words $w \in \Sigma^*$ such that $w,k \models \varphi$,
where $k$ is either $1$ or $|w|$ depending on whether the first or last position 
is regarded the \emph{output position} of $\varphi$.

\begin{example}
    The star-free language $(ab)^*$ can be defined in LTL as
    \[
        \Global( Q_a \to \next Q_b ) \wedge \Global( Q_b \wedge \next \top \to \next Q_a
        ).
    \]
    In words, at any $a$-position, the next letter is $b$. At any $b$-position
    that has a successor, the next letter is $a$.
\end{example}

\subsection{Masked unique hard-attention transformers}
Let $\Sigma$ be a finite alphabet of tokens.
A \emph{token embedding} is a function $\emb \colon \Sigma \to \Q^d$ for some $d > 0$.
A token embedding naturally extends to a homomorphism $\Sigma^* \to (\Q^d)^*$, 
where $\emb(a_1\dots a_n) = \emb(a_1)\dots\emb(a_n)$ for $a_1,\dots,a_n \in \Sigma$.

\begin{remark}
In the following we define transformers over arbitrary rational numbers since our upper bounds even hold in this setting.
We remark that all of our results also hold for the special case of fixed-precision real numbers,
i.e., with a constant number of bits for a fixed transformer regardless of the input length.
In fact, the lower bounds already hold for integers of fixed precision.
\end{remark}

\paragraph{Attention layer.}
A \emph{masked unique hard-attention} (UHA) layer of width $r > 0$ is defined by 
\begin{itemize}
\item three affine transformations $A,B \colon \Q^r \to \Q^r$ and $C \colon \Q^{2r} \to \Q^s$ with rational valued coefficients
(i.e., each is of the form $Q\bm x + \bm b$ where matrix $Q$ and vector $\bm b$ have only rational entries),
\item a mask predicate $M \colon \N \times \N \to \{\top,\bot\}$, which is defined by
$M(i,j) := \top$ (no masking), $M(i,j) := (j < i)$ (strict future masking), or $M(i,j) := (j > i)$ (strict past masking), and
\item a tie-breaking function $\tau$ selecting one element of a finite non-empty subset of $\N$, 
which is either defined as $\min$ (leftmost tie-breaking) or $\max$ (rightmost tie-breaking).
\end{itemize}
We now show how a UHA layer works on a sequence $\bm v_1,\dots,\bm v_n \in \Q^r$ with $n \ge 1$.
The \emph{score function} is defined as the dot product 
$S(\bm v_i,\bm v_j) := \langle A(\bm v_i), B(\bm v_j) \rangle$ for all $i,j \in [1,n]$.
For $i \in [1,n]$ let $U_i := \{j \in [1,n] \mid M(i,j)\}$ be the set of unmasked positions and
$B_i := \{j \in U_i \mid \forall j' \in U_i \colon S(\bm v_i,\bm v_j) \ge S(\bm v_i,\bm v_{j'})\}$
be the set of unmasked positions that maximize the score function.
We define the \emph{attention vector} at position $i \in [1,n]$ as 
$\bm a_i := \bm v_{\tau(B_i)}$ if $U_i \ne \emptyset$ and $\bm a_i := \bm 0$ otherwise.
The layer outputs the sequence $C(\bm v_1,\bm a_1),\dots,C(\bm v_n,\bm a_n)$.

\paragraph{ReLU layer.}
A ReLU layer of width $r > 0$ on input $\bm v_1,\dots,\bm v_n \in \Q^r$ 
applies for some $k \in [1,r]$ the ReLU function to the $k$-th coordinate of each $\bm v_i$, i.e.,
it outputs the sequence $\bm v'_1,\dots,\bm v'_n$ where 
$\bm v'_i := (\bm v_i[1,k-1],\max\{0,\bm v_i(k)\},\bm v_i[k+1,n])$. 
[Equivalently, one could instead allow a feed-forward network at the end of an encoder layer 
(see \citep{Hao22,barcelo2024logical}).]

\paragraph{Transformer.}
A \emph{masked unique hard-attention transformer} (UHAT) is a length-preserving function
$\mathcal{T} \colon \Sigma^* \to (\Q^s)^*$
defined by application of a token embedding followed 
by application of a fixed sequence of UHA layers and ReLU layers of matching width.
%Clearly, using an alternation of standard layers and ReLU layers, we can assume that the output of a UHAT layer 
%is an arbitrary composition of affine transformations and component-wise ReLU application.
%In particular, these compositions may use the functions $\max$ and $\min$.

\paragraph{Languages accepted by UHATs.}
To view a UHAT $\mathcal{T} \colon \Sigma^* \to (\Q^s)^*$ as a language recognizer, 
we assume that $\mathcal{T}$ is given together with an \emph{acceptance vector} $\bm t \in \Q^s$.
The recognized language $L(\mathcal{T})$ is then the set of words on which $\mathcal{T}$ outputs a sequence
$\bm v_1,\dots,\bm v_n \in \Q^s$ with $\langle \bm t,\bm v_k \rangle > 0$,
where $k \in [1,n]$ is either 1 or $n$ depending on whether the fist or last position 
is regarded the \emph{output position} of $\mathcal{T}$.

\subsection{Boolean RASP}
As an intermediate step to prove $\EXPSPACE$-hardness for UHATs,
we use Boolean RASP (\mbox{B-RASP}) as introduced by \cite{YangCA24},
who showed that B-RASP is expressively equivalent to UHATs. 
A B-RASP program is defined as follows.
Let $w = a_1 \dots a_n \in \Sigma^*$ with $n \ge 1$ be an input word.
For every $a \in \Sigma$ there is an \emph{initial} Boolean vector 
$Q_a \in \{0,1\}^n$ with $Q_a(i) = 1$ iff $a_i = a$ for all $i \in [1,n]$.
We number the initial vectors and call them $P_1,\dots,P_{|\Sigma|}$.
We now describe how vector $P_{t+1}$ can be defined from vectors $P_1,\dots,P_t$ for $t \ge |\Sigma|$.

\paragraph{Position-wise operation.} 
The vector $P_{t+1}$ can be defined by $P_{t+1}(i) := R(i)$ for some Boolean combination $R(i)$ of $\{P_1(i),\dots,P_t(i)\}$.

\paragraph{Attention operation.} 
The vector $P_{t+1}$ can be defined by either of
\begin{gather*}
P_{t+1}(i) :=\ \minatt_j [M(i,j), S(i,j)]\ V(i,j) : D(i)\\
P_{t+1}(i) :=\ \maxatt_j [M(i,j), S(i,j)]\ V(i,j) : D(i)
\end{gather*}
where
\begin{itemize}
\item $M(i,j)$ is a \emph{mask predicate} as in the definition of a UHAT,
\item $S(i,j)$ and $V(i,j)$ are Boolean combinations of $\{P_1(i),\dots,P_t(i)\} \cup \{P_1(j),\dots,P_t(j)\}$, 
called \emph{score predicate} and \emph{value predicate}, respectively,
\item $D(i)$ is a Boolean combination of $\{P_1(i),\dots,P_t(i)\}$, called \emph{default value predicate}.
\end{itemize}
The semantics of an attention operation is as follows.
For every $i \in [1,n]$, let 
\[j_i := \begin{cases}
\min\{j \in [1,n] \mid M(i,j) \text{ and } S(i,j) = 1\}, &\text{for } \minatt\\
\max\{j \in [1,n] \mid M(i,j) \text{ and } S(i,j) = 1\}, &\text{for } \maxatt.
\end{cases}\]
We then define $P_{t+1}(i) := V(i,j_i)$ if $j_i$ exists and $P_{t+1}(i) := D(i)$ otherwise.
Observe that $\minatt$ (resp.\ $\maxatt$) corresponds to leftmost (resp.\ rightmost) tie-breaking in UHATs.

A B-RASP program is given by an alphabet $\Sigma$ and a sequence of position-wise and attention operations.
We can view a B-RASP program as a language recognizer by designating one Boolean vector $Y$ as the \emph{output vector}
and either the first or last position as the \emph{output position}.
Then an input word $w = a_1 \dots a_n$ is accepted if and only if $Y(k) = 1$,
where $k$ is the output position, i.e., $k=1$ or $k=n$.

\subsection{Recurrent Neural Networks (RNN)} We use RNN as language acceptors,
as in the work of \cite{RNN-hierarchy,WGY24,WGY18}. In particular,
a \emph{Recurrent Neural Network (RNN)} $M$ can be viewed as a function 
$g: (\R^d \times \ialphabet) \to \R^d$. As in the case of transformers,
$\ialphabet$ is also mapped into $\R^k$ (for some $k$) through a token embedding
function $\emb$ (e.g. one-hot encoding) and the function $g$ actually has domain
$\R^d \times \R^k$. We have an input vector $\bar x_0 \in \R^d$ and a final 
function $f: \R^d \to \{Acc,Rej\}$, to decide whether a vector $\bar x \in \R^d$
is accepting. The semantics of acceptance is the same as that of automata, i.e.,
given $w = a_1\cdots a_n \in \ialphabet^*$, we compute $d$-vectors 
$\bar x_1, \ldots, \bar x_n$ such that $g(\bar x_i,a_{i+1}) = \bar x_{i+1}$
for each $i = 0,\ldots,n-1$. The string $w$ is \emph{accepted} by $M$ if $f(a_n)
= Acc$. 

As a computational model, it is realistic to assume RNNs with a fixed precision,
i.e., computation is always done over real numbers that can be represented with
a constant $k$ number of bits. The details of the actual representation are 
not important for our analysis. Therefore, the state-space $Q$ of the above 
RNN can be mapped to $d$-vectors over $\{0,1\}^k$ (instead of $\R$). The
following proposition is now immediate.
\begin{proposition}
    An RNN $g: (\R^d \times \ialphabet) \to \R^d$ with fixed precision $k$
    can be represented by a finite automaton with $2^{kd}$ many states.
    \label{prop:rnn-dfa}
\end{proposition}

\subsection{Size measures and succinctness}
Let $\mathcal{R}$ be a finite representation of a language, i.e., 
in our case a UHAT, LTL formula, finite automaton, RNN, or B-RASP program.
We define the size of $\mathcal{R}$, denoted by $|\mathcal{R}|$, as the length of its usual binary encoding.
In measuring succinctness of RNN, we put the precision $k$ in unary also as part of the 
size measure; since we do not want to compare a transformer that uses a
fixed precision $k$ and allow an RNN that uses a fixed precision $2^k$.
Let $\mathcal{C}_1, \mathcal{C}_2$ be classes of finite representations of languages.
We say that $\mathcal{C}_1$ can be \emph{exponentially more succinct} than $\mathcal{C}_2$ if
for every function $f \in 2^{o(n)}$ 
there is an $\mathcal{R}_1 \in \mathcal{C}_1$ such that
any $\mathcal{R}_2 \in \mathcal{C}_2$ representing the same language as $\mathcal{R}_1$ has size
$|\mathcal{R}_2| > f(|\mathcal{R}_1|)$.
Similarly, we define \emph{doubly exponentially more succinct} using functions $f \in 2^{2^{o(n)}}$ instead.
Intuitively, this means that any translation from $\mathcal{C}_1$ to $\mathcal{C}_2$ incurs, in the worst-case, 
an (doubly) exponential increase in size.
%We say that $\mathcal{C}_1$ is \emph{at least as succinct as} $\mathcal{C}_2$
%if there exists a polynomial $p$ such that for every $\mathcal{R}_1 \in \mathcal{C}_1$ 
%there is an $\mathcal{R}_2 \in \mathcal{C}_2$ representing the same language as $\mathcal{R}_1$ 
%with $|\mathcal{R}_2| \le p(|\mathcal{R}_1|)$. 
%In the case that $\mathcal{C}_1$ is at least as succinct as $\mathcal{C}_2$ and
%additionally can be (doubly) exponentially more succinct, 
%we just say that $\mathcal{C}_1$ \emph{is} (doubly) exponentially more succinct than $\mathcal{C}_2$.
